\documentclass[12pt]{article}
\usepackage[utf8]{inputenc}
\usepackage[T1]{fontenc}
\usepackage{hyperref}
\title{Yalunka Language Resources for K-12 Curriculum Development}
\author{Compiled by Claude Sonnet 5.5 (claude-sonnet-5-5) for Andr\'as Kornai}
\date{}
\begin{document}
\maketitle

\section{Introduction}
Yalunka (Dialonke, Jalonke) is a Mande language of Guinea, Sierra Leone, Mali and Senegal, closely related to Susu. The search found general and ethnographic sources and a scripture listing, but no school materials.

\section{Primary Educational Resources}
No primers or textbooks were found in the search.

\section{Linguistic and Reference Materials}
\begin{itemize}
\item Wikipedia, Yalunka people: \url{https://en.wikipedia.org/wiki/Yalunka_people}
\item University of Missouri DICE profile of Yalunka: \url{https://dice.missouri.edu/assets/docs/niger-congo/Yalunka.pdf}
\item National Library of Israel authority record, Yalunka language, which points to scholarship such as Kastenholz's work on West Mande: \url{https://nli.org.il/en/authorities/987007544479505171}
\end{itemize}

\section{Religious and Community Resources}
\begin{itemize}
\item Joshua Project, Yalunka language page. It reports a Bible translation and audio and film resources, which I did not verify: \url{https://joshuaproject.net/languages/yal}
\item Joshua Project, Yalunka (Yalun Soso) in Sierra Leone: \url{https://joshuaproject.net/index.php/people_groups/15967/SL}
\end{itemize}

\section{Digital and Audio Resources}
Audio is limited to what Joshua Project lists, which I did not check directly.

\section{Curriculum Development Context}
Yalunka is spread across four countries, and Guinea has the largest community. No school programme was found. A curriculum would need a standardized orthography and contact with Guinean and Sierra Leonean language bodies, which this search did not identify. The closeness to Susu may allow shared materials.
\end{document}